% Original PDF page 7; hash in recovery-map.json.
\recoverypage{7}
\noindent
For illustration, a normal modality can be interpreted with an indexed accessibility relation:\par

% Newly transcribed from original PDF equation (6); see recovery-map.json.
\begin{equation}
[\![\Box_i\phi]\!](w)=\forall v\,[R_i(w,v)\Rightarrow[\![\phi]\!](v)].\tag{6}\label{eq:original-6}
\end{equation}

\noindent
Different Ri and frame axioms express different profiles. A reflexive epistemic relation supports\linebreak[4]
factivity; a serial belief relation need not. A deontic ideal-world relation need not make the obligation\linebreak[4]
true in the actual world. This is an illustrative semantic recipe, not a claim that every repository\linebreak[4]
view is already compiled by Equation 6. Conditional norms, priorities, defeasible exceptions, and\linebreak[4]
contrary-to-duty obligations require their own explicit semantics; ordinary normal modal logic is not\linebreak[4]
a universal deontic solution. LogiKEy provides established semantic embeddings of deontic logics\linebreak[4]
and combinations in HOL [3].\par

\noindent
The repository's principal view families. TDFOL explicitly combines first-order, temporal, and\linebreak[4]
deontic syntax; DCEC adds agent-indexed cognitive operators and event-oriented structure. Their\linebreak[4]
source modules expose these constructors, but constructor availability is distinct from complete\linebreak[4]
backend coverage. [artifact S20--S21] The same interpreted norm may yield a deontic force/scope\linebreak[4]
view, frame-based roles, temporal constraints, event/fluent relations, graph traversal edges, and a\linebreak[4]
supported prover encoding. These are complementary task views. A predicted view distribution\linebreak[4]
estimates which interfaces deserve attention; it is not a mixture of truth values. The shared symbol\linebreak[4]
and provenance records permit joins, while each view retains its own semantic profile. [artifact\linebreak[4]
S09,S31]\linebreak[4]
Bare operator letters are unsafe identifiers: F can mean prohibition or eventually, and G can denote a\linebreak[4]
cognitive goal or a temporal always operator. Family, profile, operator kind, arguments, and nesting\linebreak[4]
must survive serialization. In particular, O(Gp) and G(Op) are not equated without a profile-specific\linebreak[4]
commutation result.\par

\subsection*{6.2  Loss-aware translation}
\noindent
A translation record must state the source/target views, profile, signature map, assumptions, supported\linebreak[4]
constructs, named losses, preserved property, and checker evidence. Preserving satisfiability is\linebreak[4]
different from preserving the truth of an arbitrary query; one-way abstraction is different from\linebreak[4]
equivalence. Translation validation examines the produced transformation rather than trusting every\linebreak[4]
invocation of a compiler [7].\linebreak[4]
A concrete legacy path makes the need visible: the inspected TDFOL-to-FOL converter removes\linebreak[4]
unary temporal and deontic operators and replaces a binary temporal formula by conjunction. This\linebreak[4]
may serve as a lossy projection, but it cannot generally carry a modal theorem. The paper does not\linebreak[4]
attribute this behavior to every adapter. For example, with p false in the actual world and true in every\linebreak[4]
ideal world, O(p) holds while p does not. On the trace p = [true, false], p holds now but Gp does\linebreak[4]
not. For strong until, q true now makes p U q true regardless of p now, contradicting the proposed\linebreak[4]
replacement p \ensuremath{\land} q. [artifact S18]\linebreak[4]
Encoding a proposition as a predicate such as obligatory(p) also does not supply its semantics. The\linebreak[4]
representation must distinguish quotation/reification from applying a predicate to a term; relational or\linebreak[4]
higher-order embeddings must interpret the modality, not simply preserve its spelling.\par

\subsection*{6.3  A conditional proof-transfer rule}
\noindent
Let \ensuremath{\Gamma} |=Ls \ensuremath{\phi} be the desired source-view consequence. A translator \ensuremath{\tau} maps the premises and goal\linebreak[4]
into Lt . Under bridge assumptions B, let RB (s, t) relate admitted source and target models. The\linebreak[4]
required directionality is:\par

% Newly transcribed from original PDF equation (7); see recovery-map.json.
\begin{equation}
s\models\Gamma\land R_B(s,t)\Rightarrow t\models\tau(\Gamma),\tag{7}\label{eq:original-7}
\end{equation}

% Newly transcribed from original PDF equation (8); see recovery-map.json.
\begin{equation}
t\models\tau(\phi)\land R_B(s,t)\Rightarrow s\models\phi.\tag{8}\label{eq:original-8}
\end{equation}

\noindent
For every admitted source model satisfying \ensuremath{\Gamma} and B, there must be a related target model; otherwise\linebreak[4]
the relation can be vacuous. If the target checker soundly establishes \ensuremath{\tau} (\ensuremath{\Gamma}) |= \ensuremath{\tau} (\ensuremath{\phi}), these conditions\linebreak[4]
transfer the consequence back. A proof sketch appears in Appendix B. This is a contract to discharge\linebreak[4]
per supported translation, not a soundness theorem about every current extractor.\linebreak[4]
Cross-source composition introduces further obligations: entity identities must align; timestamps\linebreak[4]
and units must agree; domain assumptions and exception priorities must be compatible; and the joint\par

